\scanpage{283}

\unsure{suppose} de plus \deleted{\unsure{que} \illegible{4}}\inserted{(29) \deleted{\illegible{1}} $1,\rho,\sigma$ \unsure{sont} lin.} indép. en chaque fibre i.e. $1,\rho^{-1},N_1$ lin. indép. en chaque fibre i.e. $1,\sigma^{-1},N_1$ lin. indép. en chaque fibre, \deleted{\illegible{4}}\inserted{et (30) $\operatorname{Tr}N_1\rho$ inversible.}

\margin{\illegible{2} $1,\rho^{-1},\varepsilon_1(=1+N_1)$ \illegible{1}. \illegible{2} $1,\sigma^{-1},\varepsilon_1(=1+N_1)$ \illegible{1}.}

\margin{\uline{NB} \unsure{29)} est conséquence de \unsure{30)}, qui implique \unsure{primo}, \illegible{5}$^*$
\[
\varepsilon_1=\begin{pmatrix}1&0\\1&1\end{pmatrix},\qquad
\rho=\begin{pmatrix}0&1\\a&b\end{pmatrix}\quad (a\text{ inv.}),
\]
\[
\deleted{\illegible{1}}=N_1\rho=\unsure{\cdots}=\begin{pmatrix}0&0\\0&1\end{pmatrix},
\qquad\text{donc }\sigma=\begin{pmatrix}0&-1\\-a&1-b\end{pmatrix}.
\]
$^*$ et en \unsure{normalisant} de façon à avoir $\operatorname{Tr}N_1\rho=1$.}

(NB on a $N_1\rho=-N_1\sigma=\rho+\sigma$). Ceci posé, on aura
\[
\operatorname{int}(\rho^{-1})\varepsilon_1=\operatorname{int}(\sigma^{-1})\varepsilon_1,\quad\text{soit }\varepsilon_0.
\]
(donc $\operatorname{int}(\rho)\varepsilon_0=\operatorname{int}(\sigma)\varepsilon_0=\varepsilon_1$)

\deleted{Soit $B_0$ le normalisateur de \illegible{1}\inserted{$L_0=\unsure{\text{ss-g}}$} soit $B'_0\subset B_0$ le sous-schéma en groupes défini par (74), de sorte qu'on a}
\[
\varepsilon_0=1+N_0.
\]
\deleted{\unsure{posons}} \inserted{posons}
\[
\varepsilon_0(\lambda)=1+\lambda N_0,\qquad\varepsilon_1(\lambda)=1+\lambda N_1,
\]
(sous-groupes à 1-paramètre additifs)
et soit $B_0$ le \inserted{(Borel)} normalisateur \inserted{dans $G(\mathcal M)$} des
\[
\begin{aligned}
L_0&=\text{sous-schéma en groupes }\varepsilon_0(\mathrm G_a)\\
L_1&=\text{\rule{6.3cm}{0.2pt}}\quad\varepsilon_1(\mathrm G_a),
\end{aligned}
\]
(donc $\rho(L_0)=\sigma(L_0)=L_1$),
soit $B'_0$ le sous-schéma en groupes de $B_0$ défini \inserted{en termes de $\varepsilon_0$,} par (74), de sorte que $B_0$ s'insère dans une suite exacte
\[
1\longrightarrow L_0\hookrightarrow B_0\xrightarrow{\det}\mathrm G_m\longrightarrow1,
\]
et que $\underline u\in\Gamma B_0$ implique
\[
\underline u(N_0)=\det(\underline u)N_0,\qquad\text{i.e.}
\]
\[
\underline u(\varepsilon_0(\lambda))=\varepsilon_0(\det(\underline u)\lambda).
\]

Ceci posé, considérons le morphisme
\[
\begin{aligned}
B'_0&\longrightarrow\deleted{\unsure{\mathcal X}}G(\mathcal M)\times G(\mathcal M)\times\unsure{E^1|E'}\\
\underline u&\longmapsto([\underline u,\rho],[\underline u,\sigma],\det(\underline u)).
\end{aligned}
\]
\mspage{282}

\scanpage{284}
\mspage{443}

Celui-ci prend ses valeurs dans $\mathcal X_{\rho,\sigma}$, le schéma des solutions \deleted{\unsure{des}} équations simultanées (33) à (35), et induit un iso
\[
\tag{79}B'_0\unsure{\simeq|\xrightarrow{\sim}}\mathcal X_{\rho,\sigma}.
\]

\deleted{\unsure{Cor.}} Rappelons \unsure{aussi|ainsi} pour mémoire

\uline{Corollaire} La projection
\[
\begin{aligned}
\mathcal X_\rho&\longrightarrow G(\mathcal M)\\
(\zeta,\eta,\mu)&\longmapsto\zeta
\end{aligned}
\]
est un mono, et l'image est formée des $\zeta$ satisfaisant
\[
\det\zeta=1,\quad \operatorname{Tr}\zeta\rho=\operatorname{Tr}\rho,\quad \operatorname{Tr}\zeta(\rho+\sigma)\text{ inv.}
\]
\inserted{$\unsure{-N_1\sigma=N_1\rho=\rho N_0=-\sigma N_0}$}

De $\overline{\mathrm m}$ la projection
\[
\begin{aligned}
\mathcal X_\rho&\longrightarrow G(\mathcal M)\\
(\zeta,\eta,\mu)&\longmapsto\eta
\end{aligned}
\]
est un mono, et l'image est formée des $\eta$ satisfaisant
\[
\det\eta=1,\quad\operatorname{Tr}\eta\sigma=\operatorname{Tr}\sigma,\quad\operatorname{Tr}\eta(\rho+\sigma)\text{ inv.}
\]

Enfin, on aura
\deleted{\illegible{1}$\det(\underline u)=$}
\[
\tag{80}\deleted{\unsure{\zeta}}\operatorname{Tr}\zeta(\rho+\sigma)=\operatorname{Tr}\eta(\rho+\sigma)=1/\mu.
\]

\uline{Scholie} Par transport de structure, on trouve une structure de \unsure{groupe|groupes} sur $\mathcal X_{\rho,\sigma}$.
\mspage{283}

\scanpage{285}

En l'explicitant, on trouve
\[
\begin{array}{rcl}
\text{si }\underline u&\text{correspond à}&\zeta,\eta,\mu\\
\underline u'&\text{---}&\zeta',\eta',\unsure{\mu'|\mu''}\\
\underline u''=\underline u'\underline u&\text{---}&\zeta'',\eta'',\mu''
\end{array}
\]
on aura
\[
\tag{\unsure{81|21}}\mu''=\mu\mu',\quad\zeta''=u'(\zeta)\zeta',\quad\eta''=u'(\eta)\eta'
\]
où on pose, comme d'habitude
\[
u'(\unsure{x|X})=u'\unsure{x|X}u'^{-1},
\]
\deleted{\unsure{pour}}
pour deux éléments $x,u'$ d'un groupe.

\uline{$\overline{\mathrm{III}}$}) Étude de l'équation (Suite)

\[
\tag{82}\alpha\rho(\alpha)^{-1}=\beta\sigma(\beta)^{-1}\varepsilon_1(\mu-1)
\]
\[
(\alpha,\beta\in\Gamma G(\mathcal M),\quad\mu\in\Gamma\unsure{E_S^1|\mathcal O_S^*})
\]
i.e.
\[
\tag{82bis}[\alpha,\rho]=[\beta,\sigma]\varepsilon_1(\mu-1).
\]

Soit $\mathcal G_{\rho,\sigma}$ le schéma des solutions, on a donc un morphisme
\[
\tag{83}\begin{aligned}
\mathcal G_{\rho,\sigma}&\longrightarrow\mathcal X_{\rho,\sigma}\\
(\alpha,\beta,\unsure{\mu})&\longmapsto(\alpha\rho(\alpha)^{-1},\beta\deleted{\unsure{\rho}}\inserted{\sigma}(\beta)^{-1},\mu)
\end{aligned}
\]
et il résulte du th. précédent que c'est \deleted{\illegible{1}} un épimorphisme, plus précisément,
\mspage{284}

\scanpage{286}
\mspage{\unsure{444|449}}

On a une section canonique
\[
\tag{84}
\diagram{\begin{tikzcd}[ampersand replacement=\&,column sep=large,row sep=small]
\mathcal G_{\rho,\sigma}\arrow[r]\&\mathcal X_{\rho,\sigma}\&([\underline u,\rho],[\underline u,\sigma],\det\underline u)\\
\&B'_0\arrow[ul]\arrow[u,"\wr"']\&\underline u\arrow[u,mapsto]\\
(\underline u,\underline u,\det\underline u)\&\underline u\arrow[l,mapsto]\&
\end{tikzcd}}
\]

Cette section prend ses valeurs dans la ``diagonale'', \deleted{\unsure{forme}} qui correspond aux couples $(\underline u,\mu)$ \inserted{$(\underline u\in\Gamma\mathcal M^*,\ \mu\in\Gamma\unsure{E^1|E'})$} tels que
\[
\tag{85}[\underline u,\rho]=[\underline u,\sigma]\varepsilon_1(\mu-1),
\]
qu'on \unsure{tâchera|tâche} de déterminer aussi.

Conjuguant les résultats de (I, et II) on voit que la solution générale de (82) sera \unsure{formée|fournie} \unsure{des|par},
\[
(\underline u\underline a,\underline u\underline b,\mu=\det\underline u)
\]
où $\underline u$ est une section de $B'_0$, $\underline a$ une section \deleted{locale} de $T_\rho$, $\underline b$ une section de $T_\sigma$.

De façon précise

\uline{Th.} Considérons
\[
\tag{86}\begin{aligned}
B'_0\times T_\rho\times T_\sigma&\longrightarrow G\times G\times\unsure{\mathbb E^1|E^1}\\
(\underline u,\underline a,\underline b)&\longmapsto(\underline u\underline a^{-1},\underline u\underline b^{-1},\det\underline u)
\end{aligned}
\]
\mspage{285}

\scanpage{287}

\unsure{Elle} prend ses valeurs dans $\mathcal G_{\rho,\sigma}$, et est mono --- donc induit un \uline{isomorphisme}
\[
\tag{87}B'_0\times T_\rho\times T_\sigma\xrightarrow{\sim}\mathcal G_{\rho,\sigma}
\]

De façon plus précise, soit \deleted{\unsure{localement}} $(\alpha,\beta,\mu)$ une section de $\mathcal G_{\rho,\sigma}$, alors il existe \unsure{alors} \inserted{\unsure{d'}}uniques sections \deleted{$\underline a$ de} $\underline a$ de $T_\rho$, $\underline b$ de $T_\sigma$, telles que l'on ait
\deleted{$\alpha\underline a=\beta\underline b\ (\text{soit }\underline u)$}
\[
\tag{88}\alpha\underline a,\beta\underline b\in\Gamma B'_0
\]
\deleted{$\beta\underline b$ \illegible{2}}
et on aura alors nécessairement
\[
\alpha\underline a=\beta\underline b\ (\overset{\mathrm{df}}{=}\underline u)
\]
d'où bien sûr
\[
([\alpha,\rho],[\beta,\rho])=([\underline u,\rho],[\underline u,\deleted{\unsure{\rho}}\inserted{\sigma}])
\]
et l'unique $\mu$ qui satisfasse à (82) est donc $\det\underline u$ \deleted{\illegible{1}} ; donc
\[
\tag{\unsure{89}}(\alpha,\beta,\mu)=(\underline u\underline a^{-1},\underline u\underline b^{-1},\det\underline u)\quad\text{(de façon unique)}
\]
\deleted{\unsure{Dém.}\ \unsure{Il}}

Dém. Il reste seulement à expliciter l'unicité de $\alpha,\beta$ satisfaisant à (88), ce qui est conséquence de
\[
\tag{90}T_\rho\cap B'_0=T_\sigma\cap B'_0=\{1\}
\]
qu'on avait oublié d'expliciter dans (II).
\mspage{286}

\scanpage{288}
\mspage{445}

En termes de la \inserted{seule} donnée de $\alpha$, on aura donc un $\underline a$ uniquement déterminé par la condition $\alpha\underline a\in\Gamma B'_0$, et on aura
\[
(91)\qquad \underline a=c_0\rho+d_0
\qquad
\begin{gathered}
c_0,d_0\in\Gamma(\mathcal O_S)\\
\text{déterminés par }\alpha
\end{gathered}
\]
de \unsure{$\overline{\mathrm m}$} on aura
\[
(92)\qquad \underline b=t_0\sigma+u_0
\qquad
\begin{gathered}
t_0,u_0\in\Gamma(\mathcal O_S)\\
\text{déterminés par }\beta
\end{gathered}
\]
de telle façon donc qu'on aura
\[
(93)\qquad \alpha\unsure{\underline a}=\beta\underline b
\qquad (\overset{\mathrm{dfn}}{=}\underline u)
\]

Dans le cas \unsure{numérique}
\[
\rho=\begin{pmatrix}0&1\\-1&1\end{pmatrix},\qquad
\sigma=\begin{pmatrix}0&-1\\1&0\end{pmatrix},
\]
posant
\[
\alpha=\begin{pmatrix}a&b\\c&d\end{pmatrix},\qquad
\beta=\begin{pmatrix}r&s\\t&u\end{pmatrix}
\]
on aura (cf p.\ \unsure{$398'$|$398^1$} (76)(78)) sauf erreur
\[
(94)\qquad \underline a=\frac1{\delta D}(c\rho+d)
\]
\inserted{\[
\left(=\frac1{c^2+cd+d^2}(c\rho+d)\right)
\]}
où $\delta D=c^2+cd+d^2$, --- $D$ étant le coeff \unsure{coin} de
\[
\zeta=\alpha\rho(\alpha)^{-1}=\begin{pmatrix}A&B\\C&D\end{pmatrix}
\]
\[
(95)\qquad \underline b=\frac1{\delta' U}(t\sigma+u)
\]
\inserted{\[
=\frac1{t^2+u^2}(t\sigma+u)
\]}
où $\delta' U=t^2+u^2$, --- $U$ le coeff coin de
\[
\eta=\beta\sigma(\beta)^{-1}=\begin{pmatrix}R&S\\T&\unsure{U|u}\end{pmatrix},
\]
\margin{$\delta=\det\alpha=ad-bc$\par
$\delta'=\det\beta=ru-st$}
\margin{Il faudra revenir sur ces formules dans le cas général}

le seul ennui, pour faire le lien avec les calculs ``brutaux'' de loc. cit, est dans le facteur $\tau_\infty^\alpha$ de p.\ \unsure{$398'$|$398^1$} (78), où $\alpha=\unsure{\alpha_2}(\mu)\in\mathbb Z/2.\mathbb Z$, dans le cas où $\alpha=1$ ---

\mspage{287}

\scanpage{289}

à première vue il semble y avoir incompatibilité entre la théorie (plus simple) développée ici, et ce facteur $\tau_\infty^\alpha$ --- dont \unsure{je} dois penser qu'il a de très bonnes raisons d'être là ! Et voilà --- dans loc. cit. on écrivait l'équation \unsure{des} \illegible{1} avec la possibilité d'un \uuline{signe}
\[
\zeta=\varepsilon\eta
\underbrace{\unsure{\ell_1^\nu|e_1^\nu}}_{\inserted{\substack{\varepsilon_1(\mu-1)\ \text{avec les}\\\text{notations actuelles}}}},
\qquad\varepsilon\in\{\pm1\},
\]
\margin{ici $\zeta$ et $\eta$ sont \deleted{\illegible{1}} définis comme $\alpha\rho(\alpha)^{-1}$, $\beta\sigma(\beta)^{-1}$}
c'est donc une \uline{généralisation} \unsure{p.r.} à l'équation posée au début, où on supposait $\varepsilon=+1$.
Il faudra donc \unsure{tantôt} reprendre la théorie précédente, dans le contexte plus général d'une équation
\[
(\unsure{95|96})\qquad
\underset{\inserted{\substack{\Vert\\\alpha\rho(\alpha)^{-1}}}}{\zeta}
=\varepsilon.\underset{\inserted{\substack{\Vert\\\beta\sigma(\beta)^{-1}}}}{\eta}\varepsilon_1(\mu-1)
\qquad\varepsilon\ \text{section de }\mathcal O_S
\]
où on devra avoir (en prenant les dét. des deux membres) \deleted{$\varepsilon^2=1$ \illegible{1}}
\[
\deleted{\unsure{\det}}\qquad\varepsilon^2=1
\]
--- mais il n'y a pas lieu sans doute de supposer qu'on ait $\varepsilon\in\{\pm1\}$. Mais je vais continuer d'abord les calculs dans le cas présent $\varepsilon=+1$.

\mspage{288}

\scanpage{290}
\mspage{446}

\uline{$\overline{\mathrm{IV}}$}\qquad\uline{Un épinglage canonique}

Rappelons que la donnée \inserted{(sur un schéma de base $S$)} d'une forme $G$ de $\mathrm{Gl}(2)_S$, ou $SG$ de $\mathrm{Sl}(2)_S$, \inserted{ou $PG$ de $\mathrm{GP}(1)_S$,} ou encore d'une algèbre d'\unsure{Azumay.|Azumaya} $\mathcal M$ (forme de $M_2(\mathcal O_S)$), revient à la donnée d'un fibré en droites \unsure{projectives} $\unsure{\mathcal P|P}$ sur $S$.
$\unsure{\mathcal P|P}$ peut s'interpréter comme le schéma des sous-groupes de Borel \inserted{$B$} de $G$, ou \inserted{$SB$} de $SG$, \inserted{ou $PB$ de $PG$,} ou comme le schéma des ``sous-groupes \inserted{additifs} à un paramètre ``boréliens'' '' (i.e. identiques à la partie unipotente d'un Borel convenable) \deleted{\illegible{1}}\inserted{$L$} de $G$, ou de $SG$, ou de $PG$ (NB. Les morphismes canoniques $SG\longrightarrow G\longrightarrow PG$ induisent des bijections entre \inserted{les \unsure{ens.} des} sous-groupes de Borel des trois groupes en question, et aussi entre les \unsure{ens.} des sous-groupes à 1 paramètre ``boréliens'' de ces groupes).
Également, $\unsure{\mathcal P|P}$ peut s'interpréter en termes des sections $N$ \deleted{\illegible{1}} de $\mathcal M$ qui sont non nulles sur \unsure{tte} fibre et telles que \deleted{$N^2=0$},
\inserted{\[
(1)\qquad\det N=0,\quad\operatorname{Tr}N=0\qquad\text{d'où }N^2=0
\]}
(les sections ``strictement \unsure{nilpotentes}'' \inserted{et régulières (non nulles sur chaque fibre)}), à toute telle section \unsure{correspondra|correspond} un sous groupe \inserted{à 1-par.} borélien de $SG$, comme image de $\unsure{\mathbf G_a|\mathrm G_a}$ par
\[
(2)\qquad\lambda\longmapsto\inserted{\unsure{e_N}(\lambda)=}\exp(\lambda N)=1+\lambda N.
\]
Mais il est évident que si $N$ et $N'$ sont deux telles

\mspage{289}

\scanpage{291}

sections qui sont homothétiques, i.e. $N'=aN$ avec $a\in\Gamma(S,\mathcal O_S)$, \inserted{(en fait, on aura $a\in\Gamma(S,\mathcal O_S^*)$)} alors elles définissent le \unsure{$\overline{\mathrm m}$} sous-groupe à 1 paramètre, --- donc ces dernières correspondent en fait \uline{à des} sections du fibré projectif (de dim relative 3) $\check P(\mathcal M)$, ou si on tient compte de $\operatorname{Tr}N=0$, \deleted{\illegible{1}} à des sections du fibré projectif (de dim relative 2 --- donc un fibré en plans projectifs) $\check P(\unsure{{}_1\mathcal M|1\mathcal M})$ (où $\unsure{{}_1\mathcal M|1\mathcal M}\subset\mathcal M$ est défini justement par la condition $\operatorname{Tr}N=0$) --- à savoir les sections de la \deleted{\unsure{quadrique}} ``conique'' définie par l'équation
\[
\det N=0.
\]

Les $N$ correspondent d'ailleurs aussi aux opérateurs
\[
e=e_N=\deleted{\unsure{1+N}}e_N(1)=1+N
\]
satisfaisant les conditions
\[
(3)\qquad\det e=1,\quad\operatorname{Tr}e=2
\]
ce qui équivaut en effet, en posant $e=1+N$, i.e. en définissant $N$ par
\[
N=e-1
\]
aux conditions (1) [car $\operatorname{Tr}$ et $\det$ \unsure{de deux|des} deux \unsure{opérateurs|opérations} $N$ et $e$ reliés par $e=\lambda+N$ \deleted{\unsure{se}} déterminent mutuellement par
\[
\left\{
\begin{aligned}
\operatorname{Tr}e&=2\operatorname{Tr}\lambda+\operatorname{Tr}N,&
\det e&=\lambda^2+\lambda\operatorname{Tr}N+\det N\\
\operatorname{Tr}N&=-2\operatorname{Tr}\lambda+\operatorname{Tr}e,&
\det N&=\lambda^2-\lambda\operatorname{Tr}e+\det e
\end{aligned}
\right.
\quad .
\]

\mspage{290}

\scanpage{292}
\mspage{447}

Les \deleted{opérateurs} sections de $G$ (en fait de $SG$) satisfaisant (3), \inserted{et qui ne sont $=1$ sur aucune fibre,} sont appelées les sections \deleted{\unsure{strictes|strictement}} unipotentes \inserted{régulières} --- deux telles sections \inserted{$e,e'$} sont considérées comme ``équivalentes'' si on a
\margin{terminologie de Steinberg}
\[
\begin{aligned}
(4)\qquad e'&=e^\lambda\deleted{\unsure{=e}}\inserted{\overset{\unsure{\mathrm{def}|\mathrm{dfn}}}{=}e(\lambda)}
&& (\lambda\in\Gamma\unsure{\mathbf G_a|\mathrm G_a}=\Gamma\mathcal O_S)\\
&\overset{\mathrm{def}}{=}(1+\lambda N)
&&\begin{gathered}\deleted{\text{on aura }\unsure{\text{toujours}}}\\\text{on aura }\unsure{\text{néc.}^{\mathrm t}}\quad\lambda\in\Gamma\unsure{\mathbf G_m|\mathrm G_m}\end{gathered}\\
&=1+\lambda(e-1)=(1-\lambda)+\lambda e.
\end{aligned}
\]

On voit donc que les sections de $\mathcal M$ (strictement) nilpotentes régulières correspondent \inserted{1-1} aux sections unipotentes régulières, et la relation d'équivalence d'homothétie devient celle d'équivalence par $e\longmapsto e^\lambda$ --- et les classes d'équivalence des unes ou autres correspondent aux sous-groupes additifs boréliens de $G$ (ou $SG$, ou $PG$ via $SG\longrightarrow PG$) --- ou encore $\longrightarrow$ (en prenant les normalisateurs de ceux-ci dans $G$, $SG$ ou $PG$) aux Borels de ces groupes, i.e. aux sections de $\unsure{\mathcal P|P}$ sur $S$.

Supposons donnée une telle section $\unsure{x_\infty}$ de $\unsure{\mathcal P|P}$, cela permet donc de \deleted{voir} \unsure{considérer} $\unsure{\mathcal P|P}$ comme déduit de un fibré \inserted{en droites} affines
\deleted{\illegible{3}}
\[
(5)\qquad E\quad(=\unsure{\mathcal P|P}\setminus\{\unsure{x_\infty}\}),
\]
\deleted{\unsure{Un fibré vectoriel} \illegible{4} vectoriel des translations \illegible{3} \unsure{inversible} $D$}

\mspage{291}

\scanpage{293}

et $x_\infty$ est donné également par un sous-groupe à 1-paramètre borélien
\[
L\subset SG\subset G \tag{6}
\]
(qui, par $SG\longrightarrow PG$, donne encore un sous groupe \unsure{additif} borélien $L\longrightarrow PG$ dans $PG\ldots$). Localement, $L$ sera engendré par une section unipotente régulière
\[
e=1+N\qquad (N\text{ \unsure{nilpotent} régulier}) \tag{7}
\]
comme le \unsure{s-groupe} des
\[
e^\lambda=\unsure{\underline e(\lambda)|e_{\lambda}(\lambda)}=e_N(\lambda)=1+\lambda N, \tag{8}
\]
mais comme on a dit, $e,N$ ne sont pas déterminés \unsure{intrinsèquement} en termes de $L\subset SG$. De façon précise, $L$ est muni canoniquement d'une structure de fibré vectoriel (\unsure{respectée} par son normalisateur dans $G$), \deleted{\illegible{1}} \unsure{à|:} savoir justement celle déduite par transport de structure de celle de $G_a$ par
\[
\underline e:G_a\xrightarrow{\sim}L\qquad
\underline e(\lambda)=e^\lambda=1+\lambda N,
\]
et \deleted{le choix de} $L$ s'identifie ainsi au fibré vectoriel \unsure{canonique} associé au fibré en droites $\mathcal N\subset\unsure{{}_1M}$ \deleted{\unsure{sous}} défini par les $N$
\[
L\simeq\check V(\mathcal N), \tag{9}
\]
\mspage{292}

\scanpage{294}
\mspage{448}

le choix de $e$, ou de $N$, correspond simplement à celui d'une \uuline{base} \inserted{$N$} de $\mathcal N$, ou d'une section de $L^*$ \deleted{\unsure{partout} $\unsure{\ne}$} \inserted{i.e. d'une section $e$ de $L$} qui soit $\ne1$ sur \unsure{tte} fibre.

\margin{Notons que $L$ est aussi \unsure{cano.} \unsure{isom} au fibré vectoriel des translations du fibré en droites affines $\unsure{\mathbb P|P}\setminus\{\unsure{x|x_\infty}\}$.}

Revenons à l'équation (\unsure{82|32}), et notons que $\rho$ et $\sigma$ n'y interviennent que modulo \unsure{transformations} $\rho\mapsto\lambda\rho$, $\sigma\mapsto\lambda'\sigma$ \inserted{($\lambda,\lambda'\in\Gamma G_m$)} (qui ne changent pas les commutateurs $[\alpha,\rho]$ et $[\beta,\sigma]$); \deleted{d'autre part, \unsure{les données}} \inserted{On peut l'exprimer en disant que $\rho,\sigma$ n'interviennent que par} \deleted{dans un sous-groupe additif borélien $L_1$ de $G$} leurs images dans $\Gamma(PG)$, et d'ailleurs on a
\[
\begin{aligned}
PG&\xrightarrow{\sim}\underline{\mathrm{Aut}}_{\unsure{Sgr|S\text{-gr}}}(G)
\simeq\underline{\mathrm{Aut}}_{S\text{-gr}}(SG)
\simeq\underline{\mathrm{Aut}}_{\unsure{Sgr|Sg}}(PG)\\
&\simeq\underline{\mathrm{Aut}}_{S\text{-Alg}}(M)
\simeq\underline{\mathrm{Aut}}_S(\unsure{\mathbb P|P})
\end{aligned} \tag{\unsure{10|40}}
\]
et notre remarque revient à dire que $\rho$ et $\sigma$ n'interviennent que par leur action sur $G$ (ou sur $SG$), ce qui est clair aussi sur les formules
\[
[\alpha,\rho]=\alpha\rho(\alpha)^{-1},\qquad
[\beta,\sigma]=\beta\sigma(\beta)^{-1}.
\]

D'autre part, le facteur $\varepsilon_1(\mu-1)$ de (\unsure{82|32}) peut être considéré comme une section d'un \deleted{\unsure{faux fibré}} sous-groupe additif $L_1$ de $G$ (en fait, de $SG$), qu'on suppose bien
\mspage{293}

\scanpage{295}

entendu borélien --- l'équation (\unsure{82|32}) s'écrit alors sous la forme (où $\mu$ n'intervient plus, seulement $\alpha$ et $\beta$)
\[
[\beta,\sigma]^{-1}[\alpha,\rho]\in\Gamma L_1 \tag{11}
\]
ou
\[
[\deleted{\unsure{\beta}}\inserted{\alpha},\rho]=[\beta,\sigma]\cdot\gamma
\qquad\text{avec }\gamma\in\Gamma L_1. \tag{12}
\]
Sous la forme (12), on peut la regarder comme une équation reliant $\alpha,\beta,\gamma$.

Considérons l'image \deleted{$\unsure{\widetilde L_1|\widetilde L_0}$}
\[
\widetilde L_1\subset PG\deleted{\unsure{\ell}} \tag{13}
\]
de $L_1$ par $G\longrightarrow PG$, la relation (40 bis) s'écrit alors sous la forme
\[
\widetilde\varepsilon_1\overset{\unsure{\mathrm{dfn}}}{=}
\widetilde\rho\widetilde\sigma^{-1}\in\Gamma\widetilde L_1
\qquad(\xleftarrow[\unsure{\mathrm{can}}]{\sim}\Gamma L_1). \tag{14}
\]
\deleted{Je dis que \illegible{1} la relation}
Cette section provient donc d'une section \uline{bien déterminée} \deleted{\unsure{et}} en termes de $\widetilde\rho,\widetilde\sigma$
\[
\varepsilon_1\in\Gamma L_1, \tag{15}
\]
\inserted{décrite en termes de $\rho,\sigma$ directement} par la condition qu'on ait une relation de la forme
\[
\varepsilon_1=-\lambda_0\rho\sigma^{-1}
\qquad(\lambda_0\in\Gamma G_m) \tag*{(16) = (40 bis)}
\]
\inserted{\uline{par commodité}} --- (à cause de l'application au cas particulier qui nous intéresse le plus \ldots).

\margin{\uline{NB} $\lambda_0$ est uniquement déterminé par $\rho,\sigma$ si $2$ inv. sur $S$; plus \unsure{gén\textsuperscript{t}}, si $2$ \unsure{élément} régulier de $\Gamma(S,\mathcal O_S)$. (\uline{Mais} \uline{\unsure{en car. 2}} \uline{\illegible{5}} \uline{\illegible{4}} \uline{$\unsure{\mu_2}$} \unsure{!!!}).}
\mspage{294}

\scanpage{296}
\mspage{449}

Revenons sur la situation générale d'une forme $G$ de $\unsure{\mathrm{Gl}(2)_S|\mathrm{GL}(2)_S}$, et supposons qu'elle puisse s'écrire sous la forme
\[
G=\underline{\mathrm{Aut}}_{\mathcal O_S\text{-Mod}}(F), \tag{17}
\]
où $F$ est un $\mathcal O_S$-Module \unsure{local\textsuperscript{t}} libre de rg 2,

C'est le cas dès qu'on a une section $x$ de $\unsure{\mathbb P|P}$ sur $S$, i.e. un sous-groupe additif borélien $L$ de $G$, ou de $SG$, (en fait on peut prendre alors pour $F$ l'extension de $\mathcal O_S$ par $\mathcal N$ qui détermine le fibré en droites affines $\unsure{\mathbb P|P}\setminus\{x\}$). Mais \unsure{la|le} $\unsure{\mathbb P|P}$ se réalise comme \deleted{\illegible{1}}
\[
\unsure{\mathbb P|P}=\unsure{\mathbf P|P}(F)\simeq\unsure{\check{\mathbf P}|\check P}(F) \tag{18}
\]
le fibré projectif associé à $F$, \deleted{\unsure{Ainsi}} \inserted{\unsure{i.e.}} les sections de $\unsure{\mathbb P|P}$ s'identifient aux ``droites'' $D\subset F$.

On aura
\[
M=\underline{\mathrm{End}}_{\mathcal O_S\text{-mod}}(F)\deleted{\unsure{\simeq}}\simeq\check F\otimes_{\mathcal O_S}F \tag{19}
\]
et le sous-\deleted{fibré}\inserted{Module} \deleted{vectoriel}\inserted{$\mathcal N$} associé à la section $x_D$ de $\unsure{\mathbb P|P}$ \deleted{\unsure{est}} n'est autre que \deleted{\illegible{1}}
\[
\mathcal N=D^\circ\otimes D
\qquad\bigl(D^\circ=(F/D)^\vee\subset\check F,\text{ l'orthogonal à }D\text{ dans }\check F\bigr) \tag{20}
\]
ses sections sont les $u:F\longrightarrow F$ telles que
\[
\deleted{u(F)\subset D,\quad u(D)=0,}
\]
\[
u(F)\subset D,\quad u(D)=\{0\}. \tag{21}
\]
\mspage{295}

\scanpage{297}

\uline{Proposition} \deleted{Sous les conditions précédentes} \inserted{Soit $G$ forme de $\unsure{\mathrm{Gl}(2)_S|\mathrm{GL}(2)_S}$, $\rho,\sigma$ deux sections de $G$ telles que $\widetilde\rho\widetilde\sigma^{-1}=\widetilde\varepsilon_1$ soit section \inserted{\illegible{1}} \unsure{unipot.}-\deleted{\unsure{régulière}} de $PG$, \unsure{liée} \unsure{par la} relation (16). Alors} les conditions suivantes sont équivalentes

(i) $\varepsilon_1$ est \deleted{\illegible{2}} \inserted{section} régulière (i.e. \deleted{\illegible{1}} \inserted{$\ne1$} sur chaque fibre), i.e. $\varepsilon_1=1+N_1$, avec $N_1$ une \uline{base} de $\mathcal N_1$) et $\widetilde\rho^{-1}(x_1)$ est une section de $\unsure{\mathbb P|P}$ \deleted{partout} $\ne x_1$ dans chaque fibre.

\margin{\uline{NB} On a désigné par $x_1$ la section de $\unsure{\mathbb P|P}$ sur $S$ correspondant à $L_1$.}

(ii) \unsure{Itou|Item}, avec $\widetilde\sigma^{-1}$ au lieu de $\widetilde\rho$.

\margin{$\begin{array}{l}(i\text{ bis})\\(ii\text{ bis})\end{array}\ \big|$ \unsure{Itou|Item}, avec $\widetilde\sigma,\widetilde\rho$ au lieu de $\widetilde\sigma^{-1},\widetilde\rho^{-1}$.}

(iii) On a (42), i.e. $c=\lambda_0\mathrm{Tr}\,\rho+\mathrm{Tr}\,\sigma$ est \unsure{inversible}.

\uline{Dém.} On a, grâce au fait que $\widetilde\rho\widetilde\sigma^{-1}\in\Gamma\widetilde L_1$, normalise $\widetilde L_1$, \inserted{\unsure{ou ce qui revient au} $\overline{\mathrm m}$, $L_1$,} la relation
\[
\widetilde\rho^{-1}(L_1)=\widetilde\sigma^{-1}(L_1)
\]
ou ce qui revient au $\overline{\mathrm m}$, on a
\[
\widetilde\rho^{-1}(x_1)=\deleted{\unsure{\widetilde\rho^{-1}}}\inserted{\widetilde\sigma^{-1}}(x_1), \tag{\unsure{22|24}}
\]
ce qui prouve l'équivalence de (i) et (ii).

Montrons que $(i)\Longleftrightarrow(iii)$. OPS que $S$ est le spectre d'un corps alg clos, si ça nous arrange, et OPS. $\lambda_0=1$ (quitte \deleted{\illegible{1}} à remplacer $\rho$ par $\lambda_0\rho$). Si on avait $\varepsilon_1=1$ i.e. \deleted{$\rho+\sigma$} $1=-\rho\sigma^{-1}$ i.e. $\rho+\sigma=0$, on aurait \deleted{$\unsure{\mathrm{Tr}}$} $\mathrm{Tr}(\rho+\sigma)=\mathrm{Tr}\,\rho+\mathrm{Tr}\,\sigma=0$, donc (iii) implique la régularité de $\varepsilon_1$ i.e. de $N_1$, \deleted{\illegible{1}} \inserted{OPS} donc que $\varepsilon_1$ est régulière. Il faut montrer que si on a un tel \deleted{\unsure{op}} \inserted{\unsure{endom}} unipotent régulier $\varepsilon_1$\inserted{$=1+N_1$}, de $F$, associé à une $D\subset F$,
\mspage{296}

\scanpage{298}
\mspage{450}

et \unsure{un} \deleted{\unsure{opérateur} \illegible{1}} \inserted{\unsure{endom.}} $\sigma$ de $F$, d'où
$\rho=-\varepsilon_1\sigma$, alors \deleted{\illegible{2}}
$c=\operatorname{Tr}\sigma+\operatorname{Tr}\rho=\operatorname{Tr}\sigma-\operatorname{Tr}\varepsilon_1\sigma
=-\operatorname{Tr}N_1\sigma$ est inversible sss $D+\sigma^{-1}.D=F$
i.e. $\sigma D+D=F$, \deleted{\unsure{sss}} Prenons une base
$e_0,e_1$ de $F$ telle que $e_1$ engendre $D$, et que $N_1$ s'écrive par la matrice
$\begin{pmatrix}0&0\\1&0\end{pmatrix}$,
soit $\sigma=\begin{pmatrix}a&b\\c&d\end{pmatrix}$ on a
$N_1\sigma=\begin{pmatrix}0&0\\a&b\end{pmatrix}$,
$\operatorname{Tr}N_1\sigma=b$,
$\sigma e_{\unsure{1}}=a e_0+b e_1$, et
$D+\sigma D=D+b\mathcal O_S e_{\unsure{2|0}}$
est égal à $F=\mathcal O_S^2$ sss $b$ inv., cqfd.

Supposons vérifiées ces conditions équivalentes, \deleted{\unsure{dans} $\unsure{P}$} qui s'expriment aussi en disant que $L_1$ et
\[
L_0=\widetilde\rho^{-1}(L_1)=\unsure{\widetilde\sigma}^{-1}(L_1)
\tag{23}
\]
sont distincts sur chaque fibre, i.e. \unsure{définissent} des Borels $B_1,B_0\subset G$ tels que
\[
T=B_1\cap B_0
\tag{24}
\]
soit un tore (maximal) de $G$, \deleted{\unsure{qui}}
Ce dernier est aussi la \deleted{\illegible{1}} intersection des stabilisateurs $B_0,B_1$ de $G$ des \inserted{sections}
$x_0,x_1$ de $\unsure{\mathcal P|P}$ qui correspondent à $L_0,L_1$
(satisfaisant
\[
x_1=\widetilde\rho x_0=\widetilde\sigma x_0
\tag{25}
\]
), et est \unsure{l'image} \unsure{inverse} \unsure{du s.-groupe|des s.-groupes} (\unsure{privilégié|privilégiés} \inserted{\unsure{torique|toriques}}

\mspage{297}

\scanpage{299}

\unsure{dit|dont} $\widetilde T$ dans $PG=\underline{\operatorname{Aut}}_S(\unsure{\mathcal P|P})$.
On peut dire aussi que la donnée de deux sections \deleted{\unsure{points}} distinctes
$x_0,x_1$ du fibré projectif $\unsure{\mathcal P|P}$ \deleted{\unsure{permet de}} revient à celle du Borel
$B_1$ correspondant au choix de $x_1$), \uline{et} du tore maximal $T$ de $B_1$
(qui s'écrit alors
\[
B_1=T.L_1 \qquad (\unsure{produit}\ \tfrac12\ \text{direct})
\tag{26}
\]

\margin{le couple $(B_1,T)$ s'appelle un \uline{couple de \unsure{Killing}}.}

qui permet \unsure{alors} de retrouver $x_0$ comme \unsure{un}
\inserted{correspondant à} \unsure{l'unique} Borel $B_0$ tel que
$T=B_0\cap B_1$.
(Les Borels $B_0$ tels que $B_0\ne B_1$ sur chaque fibre,
\inserted{i.e. tels que $B_0\cap B_1$ soit un tore,}
correspondent biuni-voquement aux tores maximaux $T$ de $B_1$, par la formule
\unsure{précédente} $T=B_0\cap B_1$.)

Cela revient aussi : la donnée, en plus de $B_1$, \inserted{i.e. de $x_1$,}
d'une section \inserted{$x_0$,} du fibré en droites affines
\inserted{$U_1=$}$\unsure{\mathcal P|P}\smallsetminus\{x_1\}$,
\deleted{qui \unsure{permet alors} d'identifier $U_1$}
\inserted{i.e. d'un \unsure{isom.} de $U_1$ sur le}
fibré vectoriel $L_1$ de ses translations
\[
U_1\overset{\mathrm{dfn}}{=}\unsure{\mathcal P|P}\smallsetminus\{x_1\}\xrightarrow{\sim}L_1.
\tag{27}
\]

Pour achever d'``épingler'' la situation (i.e.

\mspage{298}

\scanpage{300}
\mspage{\unsure{451|351}}

de réduire $\unsure{\mathcal P|P}$ au fibré projectif type $\mathbb P^1_S$)
il faut encore se donner une base de $L_1$ i.e. une section $x_\infty$ de
\[
L_1^*\simeq U_1\smallsetminus\{x_0\}\simeq\unsure{\mathcal P|P}\smallsetminus\{x_1,x_0\}.
\tag{28}
\]

Or la section $\varepsilon_1$ de $L_1$ est justement une section de $L_1^*$ -- on a donc un épinglage complet de la situation.
En termes de $(F,D=D_1)$ correspondant à $L_1$,

\deleted{on a
\[
F=D_1\oplus\underbrace{\sigma^{-1}(D_1)}_{D_0}
=D_1\oplus\underbrace{\rho^{-1}(D_1)}_{D_0},
\]
\unsure{OPS} (quitte à tensoriser par $D_1^{\otimes-1}$) que $D_1\simeq\mathcal O_S$,
\inserted{\unsure{on a} $N_1|D_0:$}
d'où \illegible{1}
\[
\left\{
\begin{aligned}
F&=\mathcal O_S^2\unsure{x}=\mathcal O_S e_1+\mathcal O_S e_0=\mathcal O_S e_0+\mathcal O_S e_1.\\
D_1&=\mathcal O_S e_1,\quad D_0=\mathcal O_S e_0\\
&\unsure{OPS}
\end{aligned}\right.
\tag{\unsure{29}}
\]
on aura alors $N_1=(\,\unsure{1}$}
on aura
\[
\left\{
\begin{aligned}
\rho^{-1}(D_1)&=\unsure{\sigma}^{-1}(D_1)\overset{\mathrm{dfn}}{=}D_0\\
F&=D_0\oplus D_1
\end{aligned}\right.
\tag{29}
\]
\deleted{(\unsure{30})}
\[
N_1|D_0:\quad D_0\xrightarrow{\sim}D_1
\tag{30}
\]
et (quitte à tensoriser par $D_1^{\otimes-1}$) \unsure{OPS}
$D_1\simeq\mathcal O_S$, d'où
\[
\left\{
\begin{aligned}
F&\simeq\mathcal O_S^2=\mathcal O_S e_0+\mathcal O_S e_1\\
D_0&=\mathcal O_S e_0,\quad D_1=\mathcal O_S e_1,\\
N_1e_1&=0,\quad N_1e_0=e_1,
\quad\deleted{\text{i.e. }N_1=\begin{pmatrix}0&0\\1&0\end{pmatrix}}
\end{aligned}\right.
\tag{31}
\]

\mspage{299}

\scanpage{301}

i.e.
\[
N_1=\begin{pmatrix}0&0\\1&0\end{pmatrix},\qquad
\varepsilon_1=\begin{pmatrix}1&0\\1&1\end{pmatrix}.
\tag{32}
\]

D'autre part \deleted{\unsure{on aura} $\rho D_0=\sigma\unsure{D}$}
\[
\rho D_0\unsure{=|\ne}\sigma D_0=D_1
\]
\deleted{ce qui s'exprime} \inserted{d'où}
\[
\rho\sigma^{-1}(D_1)=D_1
\]
\deleted{et on suppose $\rho,\sigma$ choisis au dessus de
$\widetilde\rho,\widetilde\sigma$ de telle façon qu'on ait
$\rho\sigma^{-1}|D_1=-1$, ce qui (compte tenu que
$\widetilde\rho\widetilde\sigma^{-1}=\widetilde\varepsilon_1$)
implique que $\rho\sigma^{-1}(F/D_1)=-1$, i.e.}
\inserted{Nous \unsure{choisissons} $\rho,\sigma$ sections en $\mathrm{Gl}(2)$)
au dessus de $\widetilde\rho,\widetilde\sigma$
(sections de $\mathrm{GP}(1)\simeq\mathrm{Gl}(2)/\mathrm G_m$
\unsure{satisfaisant} $\widetilde\rho\widetilde\sigma^{-1}=\widetilde\varepsilon_1$,}
de façon à avoir
\[
\boxed{\rho\sigma^{-1}=-\varepsilon_1}
\tag{33}
\]
ce qui détermine chacune des deux sections $\rho,\sigma$ en fonction de l'autre
(déterminée elle $\unsure{=|\overline{\mathrm m}}$ modulo un scalaire,
en \deleted{fonction} \inserted{\unsure{termes}}
de $\unsure{\widetilde\sigma}$ resp. \inserted{de} $\widetilde\rho$).
Comme
\[
\rho D_0=D_1,
\]
on aura
\[
\rho=\begin{pmatrix}0&\unsure{u|U}\\\unsure{v|V}&\unsure{w|W}\end{pmatrix},
\qquad \det\rho=-\unsure{uv|UV}\ \text{inv.}
\tag{*}
\]
et on peut lever l'\unsure{ambiguité|ambiguïté} scalaire sur $\rho$
en \deleted{\unsure{écrivant}} normalisant par
\[
\rho=\begin{pmatrix}0&1\\\unsure{v|V}&\unsure{w|W}\end{pmatrix}
\qquad \unsure{v|V}\in\Gamma\mathrm G_m.
\tag{34}
\]

\mspage{300}
